% Gauss, Werke X/1 (1917), native121–122, printed114–115.
\SourceLeaf{121}\NoSourceHead
\SourceGroup{II}
\HistoricTitle{DER GOLDENE LEHRSATZ.}
\NumberedItem{1}
\HistoricCitation{[Ein Zettel in Ea5, Kapsel 40.]}
\noindent\hspace*{\parindent}1. Es sei $p$ eine Primzahl und $p-1$ habe die Factoren $e$, $f$.
Z. B. 67; $66=6\mathpunct{.}\,11$.

2. Bekanntlich lassen sich die irrationalen Werthe von ${}^p\!\surd1$ in $e$ (6) Klassen
theilen, deren jede $f$ (11) Glieder enthält und so dass, wenn man in jeder eine Elementarwurzel
nimmt, die Exponenten der übrigen eine geometrische Reihe bilden, deren Exponenten
$={}^f\!\surd1$ Mod. $p$.

3. Die Summen dieser $e$ Klassen sind Wurzeln einer Gleichung vom Grade $e$, welche sich
leicht in jedem Falle angeben lässt.

4. \textit{Diese Gleichung ist möglich für jeden Primmodulus} $={}^f\!\surd1$ Mod $p$.
\par\noindent\hspace*{\parindent}\textls[120]{Beispiel}
\[7=\text{pr.}=1+2\mathpunct{.}\,3\]
\noindent Elementarwurzel ${}^7\!\surd1=\rho$.
\par\noindent\hspace*{\parindent}Datur aequatio pro
\[
\begin{array}{ll}
\rho+\rho^6=a&\\
\rho^2+\rho^5=b&x^3+x^2-2x-1=0\\
\rho^4+\rho^3=c&
\end{array}
\]
\noindent Diese Gleichung hat reelle Wurzeln für jeden Primmodulus $14n\pm1$.
Ex. gr. für 13 sind die Wurzeln 7, 9, 10.
\ClosingRule
\clearpage
\SourceLeaf{122}\ArithmeticHead{115}{R}{DER GOLDENE LEHRSATZ.}
\NumberedItem{2}
\HistoricCitation{[Eintragung im \textsc{Leiste}, S. 108.]}
\begin{center}\textit{Theorema generale Demonstrandum.}\end{center}
\noindent\hspace*{\parindent}1. Sit $p$ primus et si fieri potest discerpatur $p-1$ in factores $e$, $f$.
Ex. gr. $43=p$; $42=3\mathpunct{.}\,14$.

2. His positis semper datur aequatio gradus $e$, cuius radices sunt $\Sigma({}^p\!\surd1)^z$,
term. $f$, mod. $p$. Nostro casu datur aequatio
\[x^3+Axx+Bx+C=0,\]
\noindent cuius radices sunt
\[
\begin{array}{r@{\qquad}l}
1)&\varepsilon+\varepsilon^{(\rho^3)}+\varepsilon^{(\rho^6)}+\varepsilon^{(\rho^9)}+\cdots+\varepsilon^{(\rho^{39})}\\[4pt]
2)&\varepsilon^\rho+\varepsilon^{(\rho^4)}+\varepsilon^{(\rho^7)}+\cdots\\[4pt]
3)&\varepsilon^{\rho\rho}+\varepsilon^{(\rho^5)}+\cdots
\end{array}
\]
\noindent ubi $\varepsilon={}^p\!\surd1$, $\rho$ radix elem. modulo existente $p$.

Dico, hanc aequationem habere omnes radices (resid: $\int$) reales posito modulo quocunque
$m$, quando
\[m^e\equiv1\text{ modulo }p.\]
\begin{center}\textbf{BEMERKUNG.}\end{center}
\begingroup\small
\noindent\hspace*{\parindent}Vergl. die bezüglichen Bemerkungen im Artikel 17 des \textsc{Bachmann}schen
Aufsatzes »Über \textsc{Gauss}’ zahlentheoretische Arbeiten«; ferner die ebenda (in einer Fußnote
zum Artikel 15) abgedruckte Aufzeichnung, die \textsc{Gauss} in \textsc{Lamberts} \textit{Tabellen},
S. 223 eingetragen hat. In [2.] bedeutet die Abkürzung (resid: $\int$)
\textls[120]{residua summarum}, d. h. Reste der Perioden.
\par\noindent\hfill\textsc{Schlesinger}.\par
\endgroup
\ClosingRule
\clearpage
